
The hierarchical paradigm for planning and executing long-term tasks has been extensively explored over the past several decades \citep{BercherAL2019, PateriaSTQ2021}. The central idea is to decompose complex, long-horizon decision-making problems into a hierarchy of simpler subtasks, where a high-level policy operates over abstract actions, selecting subtasks rather than primitive controls, to achieve the overall objective. Each subtask can itself be formulated as a reinforcement learning problem, solved by a lower-level policy that learns to accomplish it \citep{Hengst2010}. The interaction between these hierarchical policies collectively governs the agent's overall behaviour.

By representing long-horizon problems in terms of temporally extended subtasks, hierarchical reinforcement learning and planning effectively shortens the decision horizon, a principle known as temporal abstraction \citep{BartoM2003, Dietterich2000, SuttonPS2000}. Temporal abstraction enables more efficient credit assignment across extended timescales \citep{VezhnevetsOSHJSK2017}, while decomposing tasks into subtasks typically simplifies learning and promotes more structured exploration during training \citep{NachumTLGLL2019}. These advantages make HRL a powerful framework for scaling reinforcement learning to long-horizon domains \citep{DayanH1992, VezhnevetsOSHJSK2017, NachumGLN2018}. Empirical evidence shows that HRL consistently outperforms flat RL methods across diverse settings, including continuous control \citep{FlorensaDA2017, LevyPS2018}, strategic games \citep{VezhnevetsOSHJSK2017, NachumGLN2018}, and robotic manipulation \citep{FoxKSG2017, GuptaKLLH2019}. Notably, several studies attribute these improvements primarily to enhanced exploration facilitated by subgoal-based hierarchical structures \citep{JongHST2008, NachumTLGLL2019}.


Given that large language models (LLMs) encode rich semantic knowledge about the world, such knowledge can be highly valuable for guiding embodied agents to perform high-level reasoning and planning. Motivated by this insight, \citet{AhnBBCC2022} introduced a hierarchical planning framework for robot control, in which a low-level agent executes concrete motor actions, such as opening and closing the gripper, while a high-level, LLM-based agent issues abstract commands like ``put the Coke on the counter.'' Building on this idea, subsequent studies have extensively explored the use of LLMs as task planners. These approaches can be broadly categorized into two types: comprehensive \citep{TangYTZFH2023, DalalCCS2024}, where all subgoals are planned in advance, and incremental, where subgoals are generated adaptively, allowing the planner to make dynamic adjustments based on feedback or environmental changes \citep{IchterBCKFHP2023, ZhangHL2023, WangCMLJLM2024}.


%Due to the large number of trainable parameters in LLM-based agents, optimizing them by finetuning their weights is difficult (\red{cite}). As a workaround, most of the existing methods choose to improve the LLM agents' behaviour by updating the prompting with feedback from the environments and/or let the LLM propose a list of possible high-level guidance and train a separate value function to pick them \citep{AhnBBCC2022, WangCaiCLJLM2023}.   

Due to the vast number of trainable parameters in LLM-based agents, directly optimizing them through fine-tuning is often impractical and computationally demanding \citep{HuHLKITYXYL2024,SchoeppJCYAGSZT2025}. Consequently, most existing approaches improve the behaviour of LLM agents indirectly, by refining their prompts based on environmental feedback and/or by allowing the LLM to generate multiple high-level action candidates, from which a separate value function selects the most appropriate one \citep{AhnBBCC2022, WangCaiCLJLM2023}. Beyond this common paradigm, \citet{ZhouHZZL2024,LiPSFXZ2026} propose an alternative approach that leverages the LLM as a teacher to guide a smaller student network in producing high-level decisions. Initially, the agent is trained to follow the LLM's guidance closely; however, as training proceeds, its reliance on the LLM gradually diminishes, enabling it to make independent decisions. This design removes the need for the LLM during inference, allowing the student network to be further fine-tuned in the final training stage by directly maximizing the expected return. 


%initially, the agent is trained to follow the guidance of an LLM closely, but the training evolves, 
%Combining sub-goals with these feasibility checks effectively grounds the LLM's plans in real-world constraints, helping the agent to achieve the main goal. Similarly, \citet{ZhouHZZL2024, LiPSFXZ2026} provides the agent with a set of suggested actions to execute. Initially, the agent is trained to follow the guidance of an LLM closely, but as the agent learns over time, its dependence on the LLM's suggestions decreases, allowing the agent to make independent decisions.

%\blue{
%Start from the knowledge from LLM prior.  Afterwards, fine-tuning using the pure RL way. The hierarchical approach to planning and executing long-term tasks has been extensively studied for decades \citep{BercherAL2019, PateriaSTQ2021}. 
%%\citep{PateriaSTQ2021}
%}

